\chapter{2010年天津卷（理）}

\section{单选题}

\begin{problem}
\(\i\) 是虚数单位，复数 \(\frac{-1 + 3\i}{1 + 2\i} =\)（\quad）

\choices
    {\(1 + \i\)}
    {\(5 + 5\i\)}
    {\(-5 - 5\i\)}
    {\(-1 - \i\)}
\begin{answer}
A
\end{answer}
\begin{solution}
将分子、分母同乘分母的共轭复数，得
\[
\frac{-1+3\i}{1+2\i}
=\frac{(-1+3\i)(1-2\i)}{(1+2\i)(1-2\i)}
=\frac{5+5\i}{5}=1+\i
\]
因此选 A.
\end{solution}
\end{problem}

\begin{problem}
函数 \(f(x)=2^{x}+3x\) 的零点所在的一个区间是（\quad）

\choices
    {\((-2,-1)\)}
    {\((-1,0)\)}
    {\((0,1)\)}
    {\((1,2)\)}
\begin{answer}
B
\end{answer}
\begin{solution}
函数 \(f(x)=2^x+3x\) 在 \(\R\) 上连续，且
\(f(-1)=\frac12-3=-\frac52<0\)，\(f(0)=1>0\).
由零点存在定理，\(f(x)\) 在区间 \((-1,0)\) 内至少有一个零点，因此选 B.
\end{solution}
\end{problem}

\begin{problem}
命题“若 \(f(x)\) 是奇函数，则 \(f(-x)\) 是奇函数”的否命题是（\quad）

\choices
    {若 \(f(x)\) 是偶函数，则 \(f(-x)\) 是偶函数}
    {若 \(f(x)\) 不是奇函数，则 \(f(-x)\) 不是奇函数}
    {若 \(f(-x)\) 是奇函数，则 \(f(x)\) 是奇函数}
    {若 \(f(-x)\) 不是奇函数，则 \(f(x)\) 不是奇函数}
\begin{answer}
B
\end{answer}
\begin{solution}
设命题“\(f(x)\) 是奇函数”为 \(p\)，命题“\(f(-x)\) 是奇函数”为 \(q\). 原命题是 \(p\Rightarrow q\)，其否命题为 \(\lnot p\Rightarrow\lnot q\)，即“若 \(f(x)\) 不是奇函数，则 \(f(-x)\) 不是奇函数”. 因此选 B.
\end{solution}
\end{problem}

\begin{problem}
阅读下面的程序框图，若输出 \(s\) 的值为 \(-7\)，则判断框内可填写（\quad）

\texfigure[max width=0.36\linewidth]{img_repaint/2010/tianjin_science/q4_fig1.tex}

\choices
    {\(i < 3?\)}
    {\(i < 4?\)}
    {\(i < 5?\)}
    {\(i < 6?\)}
\begin{answer}
D
\end{answer}
\begin{solution}
程序开始时 \(i=1\)，\(s=2\). 每次判断为真时，先执行 \(s\leftarrow s-i\)，再执行 \(i\leftarrow i+2\).

若判断条件为 \(i<6\)，循环中的数值依次为
\[
(i,s)=(1,2)\longmapsto(3,1)\longmapsto(5,-2)\longmapsto(7,-7)
\]
此时 \(i<6\) 不成立，输出 \(s=-7\). 而条件分别为 \(i<3\)、\(i<4\)、\(i<5\) 时，输出值分别为 \(1\)、\(-2\)、\(-2\)，均不符合题意. 因此选 D.
\end{solution}
\end{problem}

\begin{problem}
已知双曲线 \(\frac{x^2}{a^2} - \frac{y^2}{b^2} = 1\)（\(a > 0\)，\(b > 0\)）的一条渐近线方程是 \(y = \sqrt{3} x\)，它的一个焦点在抛物线 \(y^2 = 24x\) 的准线上，则双曲线的方程为（\quad）

\choices
    {\(\frac{x^2}{36} - \frac{y^2}{108} = 1\)}
    {\(\frac{x^2}{9} - \frac{y^2}{27} = 1\)}
    {\(\frac{x^2}{108} - \frac{y^2}{36} = 1\)}
    {\(\frac{x^2}{27} - \frac{y^2}{9} = 1\)}
\begin{answer}
B
\end{answer}
\begin{solution}
双曲线的渐近线为 \(y=\pm\dfrac ba x\)，故
\(\frac ba=\sqrt3\)，\(qquad b^2=3a^2\).
抛物线 \(y^2=24x=4\cdot6x\) 的准线为 \(x=-6\)，所以双曲线焦点到原点的距离满足 \(c=6\). 又
\[
c^2=a^2+b^2=4a^2
\]
从而 \(a^2=9\)，\(b^2=27\). 因此双曲线方程为
\[
\frac{x^2}{9}-\frac{y^2}{27}=1
\]
选 B.
\end{solution}
\end{problem}

\begin{problem}
已知 \(\{a_{n}\}\) 是首项为 1 的等比数列，\(S_{n}\) 是 \(\{a_{n}\}\) 的前 \(n\) 项和，且 \(9S_{3} = S_{6}\)，则数列 \(\left\{\frac{1}{a_{n}}\right\}\) 的前 5 项和为（\quad）

\choices
    {\(\frac{15}{8}\) 或 5}
    {\(\frac{31}{16}\) 或 5}
    {\(\frac{31}{16}\)}
    {\(\frac{15}{8}\)}
\begin{answer}
C
\end{answer}
\begin{solution}
设等比数列的公比为 \(q\). 因为首项为 \(1\)，所以
\(S_3=1+q+q^2\)，\(S_6=S_3(1+q^3)\).
又 \(S_3=1+q+q^2=\left(q+\frac12\right)^2+\frac34>0\). 由 \(9S_3=S_6\)，得 \(1+q^3=9\)，即 \(q=2\). 因此
\[
\sum_{n=1}^{5}\frac1{a_n}
=1+\frac12+\frac14+\frac18+\frac1{16}
=\frac{31}{16}
\]
因此选 C.
\end{solution}
\end{problem}

\begin{problem}
在 \(\triangle ABC\) 中，内角 \(A\)，\(B\)，\(C\) 的对边分别是 \(a\)，\(b\)，\(c\)，若 \(a^2 - b^2 = \sqrt{3}bc\)，\(\sin C = 2\sqrt{3}\sin B\)，则 \(A =\)（\quad）

\choices
    {\(30^{\circ}\)}
    {\(60^{\circ}\)}
    {\(120^{\circ}\)}
    {\(150^{\circ}\)}
\begin{answer}
A
\end{answer}
\begin{solution}
由正弦定理，
\(\frac cb=\frac{\sin C}{\sin B}=2\sqrt3\)，\(qquad c=2\sqrt3b\).
代入已知条件得
\(a^2-b^2=\sqrt3bc=6b^2\)，\(a^2=7b^2\).
再由余弦定理，
\[
a^2=b^2+c^2-2bc\cos A
=13b^2-4\sqrt3b^2\cos A
\]
所以 \(7=13-4\sqrt3\cos A\)，即 \(\cos A=\dfrac{\sqrt3}{2}\). 因为 \(0<A<\pi\)，故 \(A=30^\circ\)，选 A.
\end{solution}
\end{problem}

\begin{problem}
设函数 \( f(x) = \begin{cases}
\log_2 x, & x > 0,\\
\log_{\frac{1}{2}} (-x), & x < 0,
\end{cases} \) 若 \( f(a) > f(-a) \)，则实数 \( a \) 的取值范围是（\quad）

\choices
    {\((-1,0)\cup(0,1)\)}
    {\((-\infty,-1)\cup(1,+\infty)\)}
    {\((-1,0)\cup(1,+\infty)\)}
    {\((-\infty,-1)\cup(0,1)\)}
\begin{answer}
C
\end{answer}
\begin{solution}
因 \(f(a)\) 与 \(f(-a)\) 均有意义，故 \(a\ne0\).

当 \(a>0\) 时，
\(f(a)=\log_2a\)，\(qquad f(-a)=\log_{1/2}a=-\log_2a\).
不等式化为 \(\log_2a>0\)，从而 \(a>1\).

当 \(a<0\) 时，令 \(u=-a>0\)，则
\(f(a)=\log_{1/2}u=-\log_2u\)，\(qquad f(-a)=\log_2u\).
不等式化为 \(\log_2u<0\)，所以 \(0<u<1\)，即 \(-1<a<0\). 综合得
\[
a\in(-1,0)\cup(1,+\infty)
\]
选 C.
\end{solution}
\end{problem}

\begin{problem}
设集合 \(A = \{x \mid |x - a| < 1, x \in \R\}\)，\(B = \{x \mid |x - b| > 2, x \in \R\}\). 若 \(A \subset B\)，则实数 \(a\)，\(b\) 必满足（\quad）

\choices
    {\(|a + b| \leqslant 3\)}
    {\(|a + b| \geqslant 3\)}
    {\(|a - b| \leqslant 3\)}
    {\(|a - b| \geqslant 3\)}
\begin{answer}
D
\end{answer}
\begin{solution}
由集合定义，
\(A=(a-1,a+1)\)，\(B=(-\infty,b-2)\cup(b+2,+\infty)\).
区间 \(A\) 连通，若 \(A\subset B\)，则它必须完全落在 \(B\) 的某一个分支内. 因而或者
\[
a+1\le b-2
\]
得到 \(b-a\ge3\)；或者
\[
a-1\ge b+2
\]
得到 \(a-b\ge3\). 综上 \(|a-b|\ge3\)，选 D.
\end{solution}
\end{problem}

\begin{problem}
如图，用四种不同颜色给图中的 \(A\)，\(B\)，\(C\)，\(D\)，\(E\)，\(F\) 六个点涂色，要求每个点涂一种颜色，且图中每条线段的两个端点涂不同颜色，则不同的涂色方法有（\quad）

\bitmapfigure[max width=0.42\linewidth]{img/2010/tianjin_science/q10_fig1.png}

\choices
    {\(288\) 种}
    {\(264\) 种}
    {\(240\) 种}
    {\(168\) 种}
\begin{answer}
B
\end{answer}
\begin{solution}
由图可知，上方三角形为 \(\triangle ADE\)，下方三角形为 \(\triangle BCF\)，并且对应竖直线段为 \(AB\)，\(DC\)，\(EF\). 先给 \(A\)，\(D\)，\(E\) 涂色，有
\[
4\cdot3\cdot2=24
\]
种方法.

固定 \(A\)，\(D\)，\(E\) 的颜色分别为 \(1\)，\(2\)，\(3\)，记第4种颜色为 \(4\). 下面按 \(B\) 的颜色分类计数. \(B\) 不能取颜色 \(1\). 当 \(B=2\) 时，\(C\) 可取 \(1\)，\(3\)，\(4\)：若 \(C=1\)，则 \(F\) 有1种取法；若 \(C=3\)，则 \(F\) 有2种取法；若 \(C=4\)，则 \(F\) 有1种取法，共4种. 当 \(B=3\) 时，\(C\) 可取 \(1,4\)，对应的 \(F\) 各有2种取法，共4种. 当 \(B=4\) 时，\(C\) 可取 \(1,3\)，对应的 \(F\) 分别有1种、2种取法，共3种. 因而固定上方三点颜色后，下方三点有
\[
4+4+3=11
\]
种涂法. 故总涂色方法数为 \(24\cdot11=264\)，选 B.
\end{solution}
\end{problem}

\section{填空题}

\begin{problem}
甲，乙两人在 \(10\text{ 天}\) 中每天加工零件的个数用茎叶图表示如图，中间一列的数字表示零件个数的十位数，两边的数字表示零件个数个位数. 则这 \(10\text{ 天}\) 甲，乙两人日加工零件的平均数分别为\fillinblank{} 和\fillinblank{}.

\begin{center}
% 对齐的茎叶图，使用定宽列并在叶位用小空格分隔
\begin{tabular}{@{}r@{\ \ }|c@{\ \ }l@{}}
甲 & \  & 乙 \\
\hline
\(9\;8\) & \(1\) & \(9\;7\;1\) \\
\(0\;1\;3\;2\;0\) & \(2\) & \(1\;4\;2\;4\) \\
\(1\;1\;5\) & \(3\) & \(0\;2\;0\)
\end{tabular}
\end{center}
\begin{answer}
\(24\)，\(23\)
\end{answer}
\begin{solution}
从茎叶图读出甲的数据为
\(19\)，\(18\)，\(20\)，\(21\)，\(23\)，\(22\)，\(20\)，\(31\)，\(31\)，\(35\)
其和为 \(240\)，所以甲的平均数为 \(240/10=24\). 乙的数据为
\(19\)，\(17\)，\(21\)，\(24\)，\(22\)，\(24\)，\(30\)，\(32\)，\(30\)，\(11\)
其和为 \(230\)，所以乙的平均数为 \(230/10=23\). 因此两人的平均数分别为 \(24\) 和 \(23\).
\end{solution}
\end{problem}

\begin{problem}
一个几何体的三视图如图所示，则这个几何体的体积为\fillinblank{}.

\bitmapfigure[max width=0.42\linewidth]{img/2010/tianjin_science/q12_fig1.png}
\begin{answer}
\(\frac{10}{3}\)
\end{answer}
\begin{solution}
由三视图可知，该几何体由一个底面边长为 \(1\)、高为 \(2\) 的正方体柱和一个底面边长为 \(2\)、高为 \(1\) 的四棱锥组成. 因而体积为
\[
V=1^2\cdot2+\frac13\cdot2^2\cdot1
=2+\frac43=\frac{10}{3}
\]
\end{solution}
\end{problem}

\begin{problem}
已知圆 \(C\) 的圆心是直线 \(\begin{cases}
x = t,\\
y = 1 + t,
\end{cases}\)（\(t\) 为参数）与 \(x\) 轴的交点，且圆 \(C\) 与直线 \(x + y + 3 = 0\) 相切，则圆 \(C\) 的方程为\fillinblank{}.
\begin{answer}
\(\left(x+1\right)^2+y^2=2\)
\end{answer}
\begin{solution}
圆心在 \(x\) 轴上，所以由 \(y=1+t=0\) 得 \(t=-1\)，圆心为 \((-1,0)\). 圆与直线 \(x+y+3=0\) 相切，故半径为圆心到该直线的距离：
\[
r=\frac{|-1+0+3|}{\sqrt{1^2+1^2}}=\sqrt2
\]
所以圆的方程为
\[
(x+1)^2+y^2=2
\]
\end{solution}
\end{problem}

\begin{problem}
如图，四边形 \(ABCD\) 是圆 \(O\) 的内接四边形，延长 \(AB\) 和 \(DC\) 相交于点 \(P\)，若 \(\frac{PB}{PA} = \frac{1}{2}\)，\(\frac{PC}{PD} = \frac{1}{3}\)，则 \(\frac{BC}{AD}\) 的值为\fillinblank{}.

\bitmapfigure[max width=0.38\linewidth]{img/2010/tianjin_science/q14_fig1.png}
\begin{answer}
\(\frac{\sqrt{6}}{6}\)
\end{answer}
\begin{solution}
由圆内接四边形的割线定理，
\[
PA\cdot PB=PC\cdot PD
\]
设 \(PA=1\)，则 \(PB=\dfrac12\)，且由 \(PC=\dfrac{PD}{3}\) 得
\(\frac12=\frac{PD^2}{3}\)，\(qquad PD=\sqrt{\frac32}\)，\(qquad PC=\frac1{\sqrt6}\).
设 \(\angle BPC=\angle APD=\theta\). 由余弦定理，
\[
BC^2=\frac14+\frac16-\frac1{\sqrt6}\cos\theta
=\frac5{12}-\frac{\cos\theta}{\sqrt6}
\]
\[
AD^2=1+\frac32-\sqrt6\cos\theta
=\frac52-\sqrt6\cos\theta=6BC^2
\]
故
\[
\frac{BC}{AD}=\frac1{\sqrt6}=\frac{\sqrt6}{6}
\]
\end{solution}
\end{problem}

\begin{problem}
如图，在 \(\triangle ABC\) 中，\(AD \perp AB\)，\(\overrightarrow{BC} = \sqrt{3}\overrightarrow{BD}\)，\(|\overrightarrow{AD}| = 1\)，则 \(\overrightarrow{AC} \cdot \overrightarrow{AD} =\)\fillinblank{}.

\bitmapfigure[max width=0.38\linewidth]{img/2010/tianjin_science/q15_fig1.png}
\begin{answer}
\(\sqrt{3}\)
\end{answer}
\begin{solution}
由 \(D\in BC\)，有
\[
\overrightarrow{BD}=\overrightarrow{AD}-\overrightarrow{AB}
\]
因此
\[
\overrightarrow{AC}=\overrightarrow{AB}+\overrightarrow{BC}
=\overrightarrow{AB}+\sqrt3\overrightarrow{BD}
=(1-\sqrt3)\overrightarrow{AB}+\sqrt3\overrightarrow{AD}
\]
又因 \(AD\perp AB\) 且 \(|\overrightarrow{AD}|=1\)，所以
\[
\overrightarrow{AC}\cdot\overrightarrow{AD}
=(1-\sqrt3)\overrightarrow{AB}\cdot\overrightarrow{AD}
+\sqrt3\left|\overrightarrow{AD}\right|^2
=\sqrt3
\]
\end{solution}
\end{problem}

\begin{problem}
设函数 \( f(x) = x^2 - 1 \)，对任意 \( x \in \left[\frac{3}{2}, +\infty\right) \)，\( f\left(\frac{x}{m}\right) - 4m^2f(x) \leqslant f(x - 1) + 4f(m) \) 恒成立，则实数 \( m \) 的取值范围是\fillinblank{}.
\begin{answer}
\(\left(-\infty,-\frac{\sqrt{3}}{2}\right]\cup\left[\frac{\sqrt{3}}{2},+\infty\right)\)
\end{answer}
\begin{solution}
因题中含有 \(f(x/m)\)，必须有 \(m\ne0\). 将 \(f(x)=x^2-1\) 代入不等式并移项，得
\(Q_m(x):=\left(\frac1{m^2}-4m^2-1\right)x^2+2x+3\le0\)，\(x\ge\frac32\).
记 \(A_m=\dfrac1{m^2}-4m^2-1\). 若 \(A_m\ge0\)，则 \(Q_m(x)\) 在充分大的 \(x\) 处为正，不可能恒成立. 设 \(A_m<0\). 若二次函数顶点 \(-1/A_m\ge3/2\)，则其在定义区间内的最大值为
\[
Q_m\left(-\frac1{A_m}\right)=3-\frac1{A_m}>0
\]
仍不可能恒成立. 所以顶点必须在区间左侧. 此时对 \(x\ge3/2\) 有 \(Q_m'(x)=2A_mx+2<0\)，故 \(Q_m\) 在该区间上递减，只需令端点处的最大值不大于零：
\(Q_m\left(\frac32\right)=\frac94A_m+6\le0\)，\(A_m\le-\frac83\).
于是
\[
\frac1{m^2}-4m^2+\frac53\le0
\]
令 \(t=m^2>0\)，乘以正数 \(3t\) 得
\[
12t^2-5t-3\ge0
\]
该二次式的根为 \(-\dfrac13\)，\(\dfrac34\)，故 \(t\ge\dfrac34\)，即
\[
|m|\ge\frac{\sqrt3}{2}
\]
所以 \(m\) 的取值范围为
\[
\left(-\infty,-\frac{\sqrt3}{2}\right]\cup\left[\frac{\sqrt3}{2},+\infty\right)
\]
\end{solution}
\end{problem}

\section{解答题}

\begin{problem}
已知函数 \( f(x) = 2\sqrt{3}\sin x\cos x + 2\cos^2 x - 1 \)（\(x \in \R\)）.
\begin{enumerate}
\item 求函数 \( f(x) \) 的最小正周期及在区间 \( \left[0, \frac{\pi}{2}\right] \) 上的最大值和最小值；
\item 若 \(f(x_{0})=\frac{6}{5}\)，\(x_{0}\in\left[\frac{\pi}{4},\frac{\pi}{2}\right]\)，求 \( \cos 2x_{0} \) 的值.
\end{enumerate}
\end{problem}

\begin{solution}
\begin{enumerate}
\item 由二倍角公式，
\[
f(x)=\sqrt{3}\sin 2x+\cos 2x=2\sin\left(2x+\frac{\pi}{6}\right)
\]
因此函数的最小正周期为 \(\pi\).

当 \(x\in\left[0,\frac{\pi}{2}\right]\) 时，\(2x+\frac{\pi}{6}\in\left[\frac{\pi}{6},\frac{7\pi}{6}\right]\). 在此区间内，正弦函数的最大值为 \(1\)，最小值为 \(-\frac{1}{2}\)，分别在 \(2x+\frac{\pi}{6}=\frac{\pi}{2}\) 和 \(2x+\frac{\pi}{6}=\frac{7\pi}{6}\) 时取得. 所以 \(f(x)\) 在该区间上的最大值为 \(2\)，最小值为 \(-1\).
\item 由 \(f(x_0)=\frac{6}{5}\)，得
\[
\sin\left(2x_0+\frac{\pi}{6}\right)=\frac{3}{5}
\]
又因为 \(x_0\in\left[\frac{\pi}{4},\frac{\pi}{2}\right]\)，所以 \(2x_0+\frac{\pi}{6}\in\left[\frac{2\pi}{3},\frac{7\pi}{6}\right]\). 在此范围内满足上式的角位于第二象限，从而
\[
\cos\left(2x_0+\frac{\pi}{6}\right)=-\frac{4}{5}
\]
因此
\[
\cos 2x_0=\cos\left(2x_0+\frac{\pi}{6}-\frac{\pi}{6}\right)=-\frac{4}{5}\cdot\frac{\sqrt{3}}{2}+\frac{3}{5}\cdot\frac{1}{2}=\frac{3-4\sqrt{3}}{10}
\]
\end{enumerate}
\end{solution}

\begin{problem}
某射手每次射击击中目标的概率是 \(\frac{2}{3}\)，且各次射击的结果互不影响.
\begin{enumerate}
\item 假设这名射手射击 \(5\text{ 次}\)，求恰有 \(2\text{ 次}\) 击中目标的概率；
\item 假设这名射手射击 \(5\text{ 次}\)，求有 \(3\text{ 次}\) 连续击中目标，另外 \(2\text{ 次}\) 未击中目标的概率；
\item 假设这名射手射击 \(3\text{ 次}\)，每次射击，击中目标得 \(1\text{ 分}\)，未击中目标得 \(0\text{ 分}\)，在 \(3\text{ 次}\) 射击中，若有 \(2\text{ 次}\) 连续击中，而另外 \(1\text{ 次}\) 未击中，则额外加 \(1\text{ 分}\)；若 \(3\text{ 次}\) 全击中，则额外加 \(3\text{ 分}\). 记 \(\xi\) 为射手射击 \(3\text{ 次}\)后的总的分数，求 \(\xi\) 的分布列.
\end{enumerate}
\end{problem}

\begin{solution}
设一次射击击中目标和未击中目标的概率分别为 \(p=\frac{2}{3}\) 和 \(1-p=\frac{1}{3}\).
\begin{enumerate}
\item 设 \(X\) 为 \(5\) 次射击中击中目标的次数，则 \(X\) 服从二项分布. 因此
\[
P(X=2)=\mathrm C_5^2\left(\frac{2}{3}\right)^2\left(\frac{1}{3}\right)^3=\frac{40}{243}
\]
\item 恰有三次连续击中且另外两次未击中的情形为 \(HHHMM\)，\(MHHHM\)，\(MMHHH\)，其中 \(H\) 表示击中，\(M\) 表示未击中. 这三种情形互斥，所以所求概率为
\[
3\left(\frac{2}{3}\right)^3\left(\frac{1}{3}\right)^2=\frac{8}{81}
\]
\item 三次射击的结果及对应总分如下：三次均未击中时总分为 \(0\)；恰有一次击中时总分为 \(1\)；恰有两次击中且连续时总分为 \(3\)；恰有两次击中但不连续时总分为 \(2\)；三次均击中时总分为 \(6\). 因此
\(P(\xi=0)=\left(\frac{1}{3}\right)^3=\frac{1}{27}\)，
\(P(\xi=1)=\mathrm C_3^1\frac{2}{3}\left(\frac{1}{3}\right)^2=\frac{2}{9}\)，
\(P(\xi=2)=\left(\frac{2}{3}\right)^2\frac{1}{3}=\frac{4}{27}\)，
\(P(\xi=3)=2\left(\frac{2}{3}\right)^2\frac{1}{3}=\frac{8}{27}\)，且
\[
P(\xi=6)=\left(\frac{2}{3}\right)^3=\frac{8}{27}
\]
所以 \(\xi\) 的分布列为
\begin{center}
\begin{tabular}{c|ccccc}
\(\xi\) & \(0\) & \(1\) & \(2\) & \(3\) & \(6\) \\
\hline
\(P\) & \(\frac{1}{27}\) & \(\frac{2}{9}\) & \(\frac{4}{27}\) & \(\frac{8}{27}\) & \(\frac{8}{27}\)
\end{tabular}
\end{center}
\end{enumerate}
\end{solution}

\begin{problem}
如图，在长方体 \(ABCD - A_{1}B_{1}C_{1}D_{1}\) 中，\(E\)，\(F\) 分别是棱 \(BC\)，\(CC_{1}\) 上的点，\(CF = AB = 2CE\)，\(AB:AD:AA_{1} = 1:2:4\).
\begin{enumerate}
\item 求异面直线 \(EF\) 与 \(A_{1}D\) 所成角的余弦值；
\item 证明 \(AF \perp\) 平面 \(A_{1}ED\)；
\item 求二面角 \(A_{1} - ED - F\) 的正弦值.
\end{enumerate}

\bitmapfigure[max width=0.50\linewidth]{img/2010/tianjin_science/q19_fig1.png}
\end{problem}

\begin{solution}
由 \(AB:AD:AA_1=1:2:4\)，令 \(AB=1\)，则 \(AD=2\)，\(AA_1=4\). 建立空间直角坐标系，取 \(A\) 为原点，令 \(AB\)，\(AD\)，\(AA_1\) 分别为 \(x\)，\(y\)，\(z\) 轴正方向，则 \(A=(0,0,0)\)，\(B=(1,0,0)\)，\(D=(0,2,0)\)，\(C=(1,2,0)\)，\(A_1=(0,0,4)\).

由 \(CF=1\)，\(CE=\frac{1}{2}\)，得 \(E=\left(1,\frac{3}{2},0\right)\)，\(F=(1,2,1)\).

\begin{enumerate}
\item \(\overrightarrow{EF}=\left(0,\frac{1}{2},1\right)\)，\(\overrightarrow{A_1D}=(0,2,-4)\). 设异面直线所成的角为 \(\theta\)，则
\[
\cos\theta=\frac{\left|\overrightarrow{EF}\cdot\overrightarrow{A_1D}\right|}{|\overrightarrow{EF}|\,|\overrightarrow{A_1D}|}=\frac{|-3|}{\frac{\sqrt{5}}{2}\cdot2\sqrt{5}}=\frac{3}{5}
\]
\item \(\overrightarrow{AF}=(1,2,1)\)，且
\[
\overrightarrow{AF}\cdot\overrightarrow{A_1E}=(1,2,1)\cdot\left(1,\frac{3}{2},-4\right)=0
\]
\[
\overrightarrow{AF}\cdot\overrightarrow{A_1D}=(1,2,1)\cdot(0,2,-4)=0
\]
直线 \(A_1E\) 与 \(A_1D\) 是平面 \(A_1ED\) 内相交的两条直线，且平面 \(A_1ED\) 的方程为 \(x+2y+z-4=0\). 直线 \(AF\) 的方程为 \((x,y,z)=t(1,2,1)\)，它在 \(t=\frac{2}{3}\) 时与该平面相交. 因此 \(AF\perp\) 平面 \(A_1ED\).
\item 取平面 \(A_1ED\) 的法向量 \(\bs{n}_1=(1,2,1)\)，取平面 \(FED\) 的法向量
\[
\bs{n}_2=2\left(\overrightarrow{DE}\times\overrightarrow{DF}\right)=(-1,-2,1)
\]
令两平面的锐二面角为 \(\varphi\)，则
\[
|\cos\varphi|=\frac{|\bs{n}_1\cdot\bs{n}_2|}{|\bs{n}_1|\,|\bs{n}_2|}=\frac{4}{6}=\frac{2}{3}
\]
所以二面角 \(A_1-ED-F\) 的正弦值为
\[
\sin\varphi=\sqrt{1-\cos^2\varphi}=\frac{\sqrt{5}}{3}
\]
\end{enumerate}
\end{solution}

\begin{problem}
已知椭圆 \(\frac{x^2}{a^2} + \frac{y^2}{b^2} = 1 \)（\(a > b > 0\)）的离心率 \(e = \frac{\sqrt{3}}{2}\)，连接椭圆的四个顶点得到的菱形面积为 4.
\begin{enumerate}
\item 求椭圆的方程；
\item 设直线 \(l\) 与椭圆相交于不同的两点 \(A\)，\(B\). 已知点 \(A\) 的坐标为 \((-a, 0)\)，点 \(Q(0, y_{0})\) 在线段 \(AB\) 的垂直平分线上，且 \(\overrightarrow{QA} \cdot \overrightarrow{QB} = 4\). 求 \(y_{0}\) 的值.
\end{enumerate}
\end{problem}

\begin{solution}
\begin{enumerate}
\item 设椭圆的半焦距为 \(c\)，则
\(e=\frac{c}{a}=\frac{\sqrt{3}}{2}\)，\(c^2=a^2-b^2\).
从而 \(b^2=a^2-c^2=\frac{a^2}{4}\)，即 \(b=\frac{a}{2}\). 由四个顶点组成的菱形的两条对角线长分别为 \(2a\) 和 \(2b\)，所以
\[
2ab=4
\]
代入 \(b=\frac{a}{2}\)，得 \(a=2\)，\(b=1\). 因此椭圆方程为
\[
\frac{x^2}{4}+y^2=1
\]
\item 由（1）知 \(A=(-2,0)\). 过椭圆左顶点 \(A\) 的竖直直线 \(x=-2\) 是椭圆在 \(A\) 处的切线，不会与椭圆交于不同的两点，因此可设直线 \(l\) 的斜率为 \(k\)，从而 \(l\) 的方程为 \(y=k(x+2)\). 设另一交点为 \(B=(x_1,y_1)\)，将直线方程代入椭圆方程，得
\[
(1+4k^2)x^2+16k^2x+16k^2-4=0
\]
其中一个根为 \(-2\)，所以
\(x_1=\frac{2(1-4k^2)}{1+4k^2}\)，\(y_1=\frac{4k}{1+4k^2}\).
若 \(k=0\)，则 \(B=(2,0)\). 此时线段 \(AB\) 的垂直平分线为 \(y\) 轴，且
\(\overrightarrow{QA}=(-2,-y_0)\)，\(\overrightarrow{QB}=(2,-y_0)\).
由 \(\overrightarrow{QA}\cdot\overrightarrow{QB}=4\)，得 \(-4+y_0^2=4\)，故 \(y_0=\pm2\sqrt{2}\).

若 \(k\ne0\)，由 \(QA=QB\) 得
\[
4+y_0^2=x_1^2+(y_1-y_0)^2
\]
消去 \(y_0^2\) 后，代入 \(x_1\)，\(y_1\) 得
\[
2y_1y_0=x_1^2+y_1^2-4
=-\frac{48k^2}{(1+4k^2)^2}
\]
由于 \(y_1=\dfrac{4k}{1+4k^2}\ne0\)，所以
\[
y_0=-\frac{6k}{1+4k^2}
\]
又 \(\overrightarrow{QA}=(-2,-y_0)\)，\(\overrightarrow{QB}=(x_1,y_1-y_0)\)，
所以由数量积条件
\[
-2x_1-y_1y_0+y_0^2=4
\]
代入上式中的 \(x_1\)，\(y_1\)，\(y_0\)，得
\[
-\frac{4(1-4k^2)}{1+4k^2}+\frac{60k^2}{(1+4k^2)^2}=4
\]
即 \(k^2=\frac{2}{7}\). 因此
\[
y_0^2=\frac{36k^2}{(1+4k^2)^2}=\frac{56}{25}
\]
从而 \(y_0=\pm\frac{2\sqrt{14}}{5}\).

综上，\(y_0=\pm2\sqrt{2}\) 或 \(y_0=\pm\frac{2\sqrt{14}}{5}\).
\end{enumerate}
\end{solution}

\begin{problem}
已知函数 \( f(x) = x\e^{-x} \)（\(x \in \R\)）.
\begin{enumerate}
\item 求函数 \( f(x) \) 的单调区间和极值；
\item 已知函数 \( y = g(x) \) 的图象与函数 \( y = f(x) \) 的图象关于直线 \( x = 1 \) 对称. 证明：当 \( x > 1 \) 时，\( f(x) > g(x) \)；
\item 如果 \(x_{1} \neq x_{2}\)，且 \(f(x_{1}) = f(x_{2})\)，证明：\(x_{1} + x_{2} > 2\).
\end{enumerate}
\end{problem}

\begin{solution}
\begin{enumerate}
\item \(f'(x)=\e^{-x}(1-x)\). 因为 \(\e^{-x}>0\)，所以当 \(x<1\) 时 \(f'(x)>0\)，当 \(x=1\) 时 \(f'(x)=0\)，当 \(x>1\) 时 \(f'(x)<0\). 因此 \(f(x)\) 在 \((-\infty,1)\) 上单调递增，在 \((1,+\infty)\) 上单调递减，并在 \(x=1\) 处取得极大值
\[
f(1)=\e^{-1}
\]
函数没有最小值，因为当 \(x\to-\infty\) 时，\(x\e^{-x}\to-\infty\).
\item 关于直线 \(x=1\) 的对称变换把点 \((u,f(u))\) 变为点 \((2-u,f(u))\)，所以
\[
g(x)=f(2-x)=(2-x)\e^{x-2}
\]
令 \(F(x)=f(x)-g(x)\)，则 \(F(1)=0\)，且
\[
F'(x)=(1-x)\left(\e^{-x}-\e^{x-2}\right)
\]
当 \(x>1\) 时，\(1-x<0\)，且 \(-x<x-2\)，故 \(\e^{-x}<\e^{x-2}\)，从而 \(F'(x)>0\). 因此 \(F(x)>F(1)=0\)，即 \(f(x)>g(x)\).
\item 不妨设 \(x_1<x_2\). 由函数在 \((-\infty,1)\) 上递增、在 \((1,+\infty)\) 上递减，且 \(x_1\ne x_2\)，可知 \(x_1<1<x_2\). 又因为 \(f(x_2)>0\)，所以 \(x_1>0\).

若 \(x_1\leqslant2-x_2\)，则由 \(x_1<1\) 及 \(2-x_2<1\)，有
\[
f(x_1)\leqslant f(2-x_2)=g(x_2)
\]
但由（2）知 \(f(x_2)>g(x_2)\)，这与 \(f(x_1)=f(x_2)\) 矛盾. 因此 \(x_1>2-x_2\)，即 \(x_1+x_2>2\).
\end{enumerate}
\end{solution}

\begin{problem}
在数列 \(\{a_{n}\}\) 中，\(a_{1} = 0\)，且对任意 \(k \in \N^*\)，\(a_{2k-1}\)，\(a_{2k}\)，\(a_{2k+1}\) 成等差数列，其公差为 \(d_{k}\).
\begin{enumerate}
\item 若 \(d_{k} = 2k\)，证明 \(a_{2k}\)，\(a_{2k+1}\)，\(a_{2k+2}\) 成等比数列（\(k \in \N^*\)）；
\item 若对任意 \(k \in \N^*\)，\(a_{2k}\)，\(a_{2k+1}\)，\(a_{2k+2}\) 成等比数列，其公比为 \(q_k\).
\begin{enumerate}
\item 设 \(q_{1} \neq 1\). 证明 \(\left\{\frac{1}{q_{k}-1}\right\}\) 是等差数列；
\item 若 \(a_2 = 2\)，证明 \(\frac{3}{2} < 2n - \sum_{k=2}^{n} \frac{k^2}{a_k} \leqslant 2\)（\(n \geqslant 2\)）.
\end{enumerate}
\end{enumerate}
\end{problem}

\begin{solution}
\begin{enumerate}
\item 因为 \(a_{2k-1}\)，\(a_{2k}\)，\(a_{2k+1}\) 成等差数列且公差为 \(2k\)，所以
\(a_{2k}=a_{2k-1}+2k\)，\(a_{2k+1}=a_{2k}+2k\).
从而 \(a_{2k+1}=a_{2k-1}+4k\). 由 \(a_1=0\)，逐项相加得
\(a_{2k-1}=4\sum_{j=1}^{k-1}j=2k(k-1)\)，\(a_{2k}=2k^2\)，\(a_{2k+1}=2k(k+1)\).
因此
\[
a_{2k+1}^{2}=4k^2(k+1)^2=a_{2k}a_{2k+2}
\]
所以 \(a_{2k}\)，\(a_{2k+1}\)，\(a_{2k+2}\) 成等比数列.
\item
\begin{enumerate}
\item 由“公比为 \(q_k\)”的定义，\(a_{2k+1}=q_ka_{2k}\)，\(a_{2k+2}=q_ka_{2k+1}=q_k^2a_{2k}\). 又由下一组三项成等差数列，
\[
a_{2k+3}=2a_{2k+2}-a_{2k+1}=q_k(2q_k-1)a_{2k}
\]
若 \(q_k\ne0\)，则 \(a_{2k+2}=q_k^2a_{2k}\ne0\)，故下一组三项的公比可写为
\[
q_{k+1}=\frac{a_{2k+3}}{a_{2k+2}}=\frac{2q_k-1}{q_k}=2-\frac{1}{q_k}
\]
第一组三项中 \(a_1=0\)，故 \(a_3=2a_2\). 又因 \(q_1\) 是三项 \(a_2\)，\(a_3\)，\(a_4\) 的公比，按公比定义 \(a_2\ne0\)，且 \(a_3=q_1a_2\)，从而 \(q_1=2\). 于是 \(q_1\ne0\)，上面的递推式首先给出 \(q_2\). 若已知 \(q_k=\dfrac{k+1}{k}\)，则 \(q_k\ne0\)，且
\[
q_{k+1}=2-\frac{k}{k+1}=\frac{k+2}{k+1}
\]
由数学归纳法，
\(q_k=\frac{k+1}{k}\)，\((k\in\N^*)\)
故各 \(q_k\ne0\)，\(1\)，并且
\[
\frac{1}{q_{k+1}-1}=\frac{1}{1-1/q_k}=\frac{q_k}{q_k-1}=1+\frac{1}{q_k-1}
\]
所以数列 \(\left\{\frac{1}{q_k-1}\right\}\) 是公差为 \(1\) 的等差数列.
\item 由 \(a_2=2\) 和第一组三项的等差关系，得 \(a_3=4\)，故 \(q_1=2\). 由（i）可知
\(\frac{1}{q_k-1}=k\)，\(q_k=\frac{k+1}{k}\).
先有 \(a_2=2\)，且 \(a_3=q_1a_2=4\). 若 \(a_{2k}=2k^2\)，则由等比关系
\(a_{2k+1}=q_ka_{2k}=\frac{k+1}{k}\cdot2k^2=2k(k+1)\)，\(a_{2k+2}=q_k^2a_{2k}=\left(\frac{k+1}{k}\right)^2\cdot2k^2=2(k+1)^2\).
由 \(k=1\) 的初值开始归纳，得到
\(a_{2k}=2k^2\)，\(a_{2k+1}=2k(k+1)\)，\((k\geqslant1)\).
因此，对 \(k\geqslant2\)，
\[
\frac{k^2}{a_k}=\begin{cases}
2,&k\text{为偶数},\\
2+\dfrac{2}{k^2-1},&k\text{为奇数}
\end{cases}
\]
于是
\[
2n-\sum_{k=2}^{n}\frac{k^2}{a_k}=2-\sum_{\substack{3\leqslant k\leqslant n\\ k\text{为奇数}}}\frac{2}{k^2-1}
\]
当 \(n=2m\) 时，利用裂项
\[
\sum_{j=1}^{m-1}\frac{2}{(2j+1)^2-1}=\sum_{j=1}^{m-1}\left(\frac{1}{2j}-\frac{1}{2j+2}\right)=\frac{1}{2}-\frac{1}{2m}
\]
从而
\[
2n-\sum_{k=2}^{n}\frac{k^2}{a_k}=\frac{3}{2}+\frac{1}{n}
\]
当 \(n=2m+1\) 时，
\[
\sum_{j=1}^{m}\frac{2}{(2j+1)^2-1}=\sum_{j=1}^{m}\left(\frac{1}{2j}-\frac{1}{2j+2}\right)=\frac{1}{2}-\frac{1}{2m+2}
\]
从而
\[
2n-\sum_{k=2}^{n}\frac{k^2}{a_k}=\frac{3}{2}+\frac{1}{n+1}
\]
当 \(n\geqslant2\) 时，上述两种情形的结果均大于 \(\frac{3}{2}\) 且不超过 \(2\)，故不等式成立.
\end{enumerate}
\end{enumerate}
\end{solution}
